\section{Limitaciones del aprendizaje profundo}

\subsection{Revisitando la generalización: el experimento de etiquetas aleatorias}

Amini presenta un experimento clásico, del artículo de Zhang et al., publicado en 2017 en ICLR: \textit{Understanding Deep Neural Networks Requires Rethinking Generalization}.

\begin{knowledgebox}{Diseño del experimento}
Los investigadores tomaron las imágenes del conjunto de datos ImageNet junto con sus etiquetas y, para cada imagen de forma independiente, lanzaron un dado de $k$ caras ($k$ es el número total de categorías), sustituyendo la etiqueta original por el resultado aleatorio. La consecuencia de esto es que dos imágenes de la misma categoría pueden terminar asignadas a etiquetas completamente distintas, y la asociación semántica entre la etiqueta y el contenido de la imagen se pierde por completo.
\end{knowledgebox}

\begin{figure}[H]
\centering
\includegraphics[width=0.85\textwidth]{slides-images/slide-19.jpg}
\caption{Experimento de etiquetas aleatorias: las etiquetas de las imágenes se mezclan por completo\protect\footnotemark}
\end{figure}
\footnotetext{Diapositiva 19. Cita a Zhang+ ICLR 2017.}

\subsubsection{Resultados del experimento}

Los investigadores entrenaron redes neuronales profundas con distintos grados de aleatorización de las etiquetas, observando el desempeño en el conjunto de entrenamiento y en el conjunto de prueba:

\begin{figure}[H]
\centering
\includegraphics[width=0.85\textwidth]{slides-images/slide-22.jpg}
\caption{Hallazgo clave: incluso con etiquetas completamente aleatorias, la precisión de entrenamiento puede llegar al 100\%\protect\footnotemark}
\end{figure}
\footnotetext{Diapositiva 22. La línea punteada roja marca que la precisión de entrenamiento se mantiene siempre cerca del 100\%.}

\begin{importantbox}{Hallazgo central: la sorprendente capacidad de ajuste de las redes profundas}
\begin{itemize}
    \item \textbf{Precisión de prueba}: conforme aumenta el grado de aleatorización de las etiquetas, la precisión en el conjunto de prueba disminuye continuamente (algo intuitivo).
    \item \textbf{Precisión de entrenamiento}: sin importar qué tan aleatorias sean las etiquetas, la red profunda siempre puede alcanzar una precisión cercana al 100\% en el conjunto de entrenamiento.
\end{itemize}
Esto significa que las redes profundas actuales tienen capacidad suficiente para ``memorizar'' todo el conjunto de entrenamiento, incluso cuando en los datos no existe ningún patrón que se pueda aprender.
\end{importantbox}

\subsection{La doble filo del aproximador de funciones}

A partir del experimento anterior, Amini profundiza en este fenómeno desde la perspectiva del ajuste de funciones:

\begin{figure}[H]
\centering
\includegraphics[width=0.85\textwidth]{slides-images/slide-25.jpg}
\caption{La red neuronal como aproximador de funciones: un desempeño excelente en las regiones con datos de entrenamiento\protect\footnotemark}
\end{figure}
\footnotetext{Diapositiva 25.}

La red neuronal puede ajustar muy bien los valores de la función cerca de los puntos de datos de entrenamiento. Dado un nuevo punto de datos (en morado), si cae dentro del rango de la distribución de los datos de entrenamiento, la red puede dar una predicción razonable. Pero el problema central es el siguiente:

\begin{figure}[H]
\centering
\includegraphics[width=0.85\textwidth]{slides-images/slide-28.jpg}
\caption{El comportamiento fuera de la distribución de los datos es impredecible\protect\footnotemark}
\end{figure}
\footnotetext{Diapositiva 28.}

\begin{warningbox}{El problema out-of-distribution}
Cuando los datos de entrada quedan fuera de la distribución de entrenamiento (out-of-distribution, OOD), el comportamiento de la red neuronal es impredecible. En las regiones sin datos de entrenamiento, la función ajustada puede producir valores de salida arbitrarios, y la red no te ``avisa'' que no está segura. Esto plantea una pregunta central: \textbf{¿cómo sabemos qué es lo que el modelo no sabe?}
\end{warningbox}

\subsection{El aprendizaje profundo no es alquimia}

\begin{figure}[H]
\centering
\includegraphics[width=0.85\textwidth]{slides-images/slide-29.jpg}
\caption{Deep Learning $\neq$ Alchemy\protect\footnotemark}
\end{figure}
\footnotetext{Diapositiva 29. Cita a U. Muller, 6.S191 2018.}

Amini enfatiza: el aprendizaje profundo \textbf{no} es una ``caja negra mágica'' en la que puedes meter cualquier cosa y obtener una salida perfecta. El principio clásico de ``garbage in, garbage out'' (basura entra, basura sale) se aplica por completo. Antes de desplegar una solución de IA hay que considerar dos preguntas centrales:
\begin{enumerate}
    \item ¿Esta tarea realmente es adecuada para resolverse con aprendizaje profundo?
    \item ¿Puedes recopilar y depurar datos de alta calidad para esta tarea?
\end{enumerate}

\subsection{Modo de falla uno: sesgo en los datos}

\begin{figure}[H]
\centering
\begin{subfigure}[b]{0.48\textwidth}
    \includegraphics[width=\textwidth]{slides-images/slide-30.jpg}
    \caption{En las imágenes de entrenamiento, una gran cantidad de perros aparecen con la lengua afuera}
\end{subfigure}
\hfill
\begin{subfigure}[b]{0.48\textwidth}
    \includegraphics[width=\textwidth]{slides-images/slide-31.jpg}
    \caption{Zonas rosas que aparecen cuando la CNN colorea imágenes en blanco y negro}
\end{subfigure}
\caption{Falla causada por el sesgo en los datos: el modelo de coloreado de imágenes\protect\footnotemark}
\end{figure}
\footnotetext{Diapositivas 30--31. Cita a P. Isola 6.869.}

El curso presenta un ejemplo elocuente: entrenar una CNN para convertir fotografías en blanco y negro a color. Cuando en la fotografía de salida de un perro aparecía una mancha rosa en la zona de la mandíbula, la razón era que, en el conjunto de entrenamiento, una gran cantidad de fotografías de perros los mostraban con la lengua rosa afuera; el modelo aprendió la asociación estadística de que ``cerca de la mandíbula del perro debería haber rosa'', y no una verdadera comprensión del significado del color.

\subsection{Modo de falla dos: incertidumbre en escenarios críticos para la seguridad}

\begin{figure}[H]
\centering
\includegraphics[width=0.85\textwidth]{slides-images/slide-32.jpg}
\caption{Caso de accidente fatal del Autopilot de Tesla\protect\footnotemark}
\end{figure}
\footnotetext{Diapositiva 32. Cita a ABC News.}

Amini menciona un accidente fatal real de conducción autónoma: un Tesla, en el mismo tramo de carretera, se desvió repetidamente hacia la barrera divisoria hasta que finalmente terminó en una colisión. La investigación reveló que la barrera divisoria de ese tramo se había construido después de la recopilación de los datos de entrenamiento: el sistema de conducción autónoma se encontró con un escenario \textbf{fuera de la distribución de entrenamiento}.

\begin{figure}[H]
\centering
\includegraphics[width=0.85\textwidth]{slides-images/slide-33.jpg}
\caption{La incertidumbre es crucial en las aplicaciones críticas para la seguridad\protect\footnotemark}
\end{figure}
\footnotetext{Diapositiva 33.}

\begin{importantbox}{La incertidumbre en el aprendizaje profundo}
En los siguientes escenarios, detectar y cuantificar la incertidumbre de manera confiable es especialmente importante:
\begin{itemize}
    \item \textbf{Aplicaciones críticas para la seguridad}: conducción autónoma, diagnóstico médico, reconocimiento facial
    \item \textbf{Problemas de calidad de los datos}: desbalance de clases (imbalance), ruido en los datos (data noise)
    \item \textbf{Detección fuera de distribución}: cuando los datos de entrada difieren de manera significativa de la distribución de entrenamiento, el modelo debería ``saber que no sabe''
\end{itemize}
\end{importantbox}

\subsection{Modo de falla tres: ataques adversarios}

Los ataques adversarios (Adversarial Attacks) son un problema de seguridad clásico y muy importante en el aprendizaje profundo.

\begin{figure}[H]
\centering
\includegraphics[width=0.85\textwidth]{slides-images/slide-35.jpg}
\caption{Ejemplo de ataque adversario: Temple (97\%) $\rightarrow$ Ostrich (98\%)\protect\footnotemark}
\end{figure}
\footnotetext{Diapositiva 35.}

\subsubsection{El fundamento matemático de los ataques adversarios}

Recordemos la fórmula central del descenso de gradiente:

$$
W \leftarrow W - \eta \frac{\partial J(W, x, y)}{\partial W}
$$

donde $W$ son los pesos, $\eta$ es la tasa de aprendizaje, $J$ es la función de pérdida y $(x, y)$ son la entrada y la etiqueta fijas. Durante el entrenamiento, \textbf{fijamos la entrada $x$ y la etiqueta $y$}, y optimizamos los pesos $W$ para minimizar la pérdida.

\begin{figure}[H]
\centering
\includegraphics[width=0.85\textwidth]{slides-images/slide-40.jpg}
\caption{Fórmula para generar la imagen adversaria\protect\footnotemark}
\end{figure}
\footnotetext{Diapositiva 40.}

El ataque adversario invierte por completo este proceso:

$$
x \leftarrow x + \eta \frac{\partial J(W, x, y)}{\partial x}
$$

\begin{importantbox}{Ataque adversario vs. entrenamiento normal}
\begin{itemize}
    \item \textbf{Entrenamiento normal}: se fija la entrada $(x, y)$ y se optimizan los pesos $W$ para \textbf{minimizar} la pérdida
    \item \textbf{Ataque adversario}: se fijan los pesos $W$ y la etiqueta $y$, y se modifica la entrada $x$ para \textbf{maximizar} la pérdida
\end{itemize}
La diferencia clave: el objeto respecto al cual se calcula el gradiente pasa de $\partial W$ a $\partial x$, y la dirección de la optimización pasa de un signo menos a un signo más.
\end{importantbox}

\subsubsection{Ejemplos adversarios en el mundo físico}

Los ataques adversarios no solo existen en el espacio digital, también pueden realizarse en el mundo físico. Amini presenta un ejemplo adversario físico fabricado por estudiantes del MIT con impresión 3D: una tortuga impresa en 3D con una textura cuidadosamente diseñada que, desde diversos ángulos y bajo distintas condiciones de iluminación, siempre es clasificada por el clasificador como ``rifle'' en lugar de ``tortuga''.

\begin{figure}[H]
\centering
\includegraphics[width=0.85\textwidth]{slides-images/slide-42.jpg}
\caption{Ejemplo adversario físico impreso en 3D: una tortuga identificada como rifle\protect\footnotemark}
\end{figure}
\footnotetext{Diapositiva 42. Cita a Athalye+ ICML 2018.}

\begin{warningbox}{La amenaza real de los ataques adversarios}
La existencia de ejemplos adversarios implica que, en escenarios críticos para la seguridad (como el reconocimiento de señales de tránsito en la conducción autónoma o el reconocimiento facial en sistemas de seguridad), una perturbación mínima y cuidadosamente diseñada puede provocar que el sistema emita un juicio completamente erróneo. Esto no es solo un problema académico, sino un riesgo de seguridad que debe considerarse necesariamente al desplegar sistemas de IA en la práctica.
\end{warningbox}

\subsection{Síntesis de las limitaciones}

\begin{figure}[H]
\centering
\includegraphics[width=0.85\textwidth]{slides-images/slide-46.jpg}
\caption{Lista completa de las limitaciones de las redes neuronales\protect\footnotemark}
\end{figure}
\footnotetext{Diapositiva 46.}

El curso resume las principales limitaciones de las redes neuronales:

\begin{table}[H]
\centering
\begin{tabular}{ll}
\toprule
\textbf{Limitación} & \textbf{Descripción} \\
\midrule
Hambrienta de datos (Data hungry) & Por lo general requiere millones de muestras de entrenamiento \\
Intensiva en cómputo (Computationally intensive) & Tanto el entrenamiento como el despliegue requieren GPU \\
Vulnerable a ataques adversarios (Adversarial examples) & Fácil de engañar con perturbaciones cuidadosamente diseñadas \\
Sesgo algorítmico (Algorithmic bias) & Puede amplificar los sesgos sociales presentes en los datos \\
Mala representación de la incertidumbre (Uncertainty) & Es difícil saber cuándo el modelo ``no sabe'' \\
Caja negra inexplicable (Black boxes) & Difícil de entender y de confiar en ella \\
Requiere conocimiento experto (Expert knowledge) & El diseño de la arquitectura y el ajuste de hiperparámetros dependen de la experiencia \\
\textbf{Dificultad para codificar estructura} (Encode structure) & Difícil incorporar conocimiento previo al aprendizaje \\
\textbf{Dificultad de extrapolación} (Extrapolation) & Difícil ir más allá del rango de la distribución de los datos de entrenamiento \\
\bottomrule
\end{tabular}
\caption{Principales limitaciones de las redes neuronales}
\end{table}

\begin{knowledgebox}{Las limitaciones engendran nuevas fronteras}
Amini señala: como técnicos y científicos, las limitaciones representan precisamente oportunidades de desarrollo. Los Diffusion Models y los LLM que se presentan en la segunda mitad del curso son justamente una respuesta a estos dos retos centrales: la ``dificultad de extrapolación'' y la ``dificultad para codificar estructura''; al modelar la propia distribución de los datos, se espera que los modelos base (Foundation Models) logren una capacidad de generalización más sólida.
\end{knowledgebox}

\subsection{Resumen de la sección}

Las limitaciones del aprendizaje profundo abarcan múltiples aspectos: una capacidad de generalización insuficiente (el experimento de etiquetas aleatorias reveló una sorprendente capacidad de sobreajuste), errores sistemáticos causados por el sesgo en los datos, retos de incertidumbre en escenarios críticos para la seguridad, y la amenaza que representan los ataques adversarios para la robustez del modelo. Comprender estas limitaciones es un requisito para usar la tecnología de IA de manera responsable, y también es un motor que impulsa el desarrollo de la próxima generación de algoritmos.

%% =====================================================
